\section{Worked examples}
\label{sec:examples}

\subsection{A demand-passing arithmetic pipeline}

Consider a sequence of approximate operations.  Stage \(n+1\) produces an
output at precision \(r\) from a stage-\(n\) input at precision \(r+c_n\), where
\(c_n\) is a fixed primitive-recursive loss.  Thus
\(D(n,r)=r+c_n\).  The backward schedule is
\[
 h_{N,k}(t)=k+\sum_{j=t}^{N-1}c_j.
\]
This identity is proved by induction from the schedule equations.

Suppose \(P(n,r)\) says that the stage-\(n\) rational approximation has error at
most \(2^{-r}\), and the local arithmetic lemma proves
\[
 P(n,r+c_n)\to P(n+1,r).
\]
The demand theorem gives \(P(N,k)\) from one base approximation at
\(k+\sum_{j<N}c_j\).  No global monotonicity theorem for the entire pipeline is
needed.  In Lean, the sum is a \(\PR\) accumulator, the local estimate is an
object derivation, and \code{primitiveRecursionAt} carries the invariant.

\subsection{Dead-end removal in a finite approximation tree}

Let level zero have candidates \(0,1\), level one have \(0,1\), and level two
have only candidate \(1\).  All candidates are locally good.  Compatibility
edges are
\[
 0\to0,\qquad 1\to1
\]
from level zero to one, and only
\[
 1\to1
\]
from level one to two.  Choosing the least good candidate at every level starts
with zero and becomes stuck.  Backward viability at horizon two marks level-one
candidate zero nonviable and propagates this information to level zero.  The
least viable path is \((1,1,1)\).

This small example captures a general issue in chart selection, root-cluster
matching, and finite descent: local admissibility is not extendibility.  The
viability recurrence is the exact finite correction.

\subsection{Independent square-root intervals}

Fix a positive rational \(q\).  At precision \(n\), a certificate is a dyadic
interval \([L,U]\) such that
\[
 0\le L\le U,
 \qquad L^2\le q\le U^2,
 \qquad U-L\le2^{-n-2}.
\]
The observation is the midpoint \((L+U)/2\).  Any two valid intervals at
precisions \(n,m\) both contain \(\sqrt q\) in the ordered-field semantics, so
their midpoints differ by at most
\[
 2^{-n-2}+2^{-m-2}.
\]
The same inequality can be proved without mentioning a completed real: from
\(L^2\le q\le U^2\) and nonnegativity, two valid intervals must overlap, for
otherwise the upper endpoint of one would lie strictly below the lower endpoint
of the other and squaring would contradict the common bracket.  Overlap plus
the two widths gives the midpoint bound.

Therefore arbitrary independently produced valid intervals yield a regular
name by \cref{thm:stable-erasure}.  No nested-interval producer is required.
This is an instructive improvement over a path-based construction: nesting is
sufficient but not necessary.

A verifier checks all displayed rational inequalities by cross multiplication.
A producer may use bisection.  The object theorem requires only the verifier's
soundness and pairwise interval comparison.

\subsection{Finite-consumer metastability}

Suppose a bounded routine evaluates a rational sequence at at most eight
indices, each between a base \(N\) and \(N+3N+20\), and proves a conclusion
provided all queried values differ by at most \(2^{-r}\).  Set
\(G(N)=3N+20\).  A metastability certificate for this \(G\) supplies a stable
window containing every query.  The routine's trace then proves its conclusion.

Nothing follows about indices outside that window, and no Cauchy modulus has
been constructed.  Nevertheless the finite downstream theorem is fully valid.
This is the intended use of metastability in iPRA.

\subsection{Height bound and complete enumeration}

Let \(P(x)\) be a decidable Diophantine predicate and suppose all codes with
height at most \(H\) lie below a finite code bound \(C(H)\).  Given a height
certificate, scan \(x<C(H)\), retain exactly those satisfying \(P\), and remove
duplicates.  The resulting list is complete.  Given a complete list, compute
the maximum height.  Both directions are finite traces with straightforward
Boolean invariants.

The example shows exactly where a Faltings- or Siegel-type application remains
difficult: not in enumerating or checking points once \(H\) is known, but in
producing the uniform height certificate.
